% Part: lambda-calculus
% Chapter: introduction
% Section: currying

\documentclass[../../../include/open-logic-section]{subfiles}

\begin{document}

\olfileid{lam}{int}{cur}
\olsection{કરિંગ}
\sourcecorrection{OLLAM-001}{આ વિભાગ પ્રસ્તાવના પ્રકરણમાં
આયાત થાય છે, પરંતુ સ્થિર અંગ્રેજી મૂળના ટિપ્પણી-મથાળામાં
\texttt{representability} અને ફાઇલ-ઓળખમાં~\texttt{rep} છે.
પ્રકરણની વાસ્તવિક જગ્યા પ્રમાણે અહીં~\texttt{int} મૂક્યું છે.}

$\lambd$-અમૂર્ત અભિવ્યક્તિ~$\lambd[x][M]$ એક દલીલવાળું
વિધેય દર્શાવે છે. અનેક દલીલો સ્વીકારતું વિધેય વ્યાખ્યાયિત કરવું
હોય ત્યારે આ મર્યાદા લાગે છે. એક રસ્તો એ છે કે~$\lambd$-કલન
વિસ્તારીને જોડી, ત્રિક વગેરે બનાવવાની છૂટ આપીએ; પછી, દાખલા
તરીકે, ત્રિસ્થાની વિધેય~$\lambd[x][M]$ની દલીલ ત્રિક હશે.
પરંતુ આ માટે \emph{કરિંગ} વધુ અનુકૂળ છે.

દાખલો જોઈએ. થોડા સમય માટે માની લઈએ કે~$\lambd$-કલનમાં
$+$ ક્રિયા છે. સરવાળાનું વિધેય~$2$-સ્થાની છે, એટલે બે દલીલો
સ્વીકારે છે. પરંતુ~$\lambd$-અમૂર્તીકરણ માત્ર એક દલીલવાળાં વિધેયો
આપે છે: વાક્યરચનામાં $\lambd[(x,y)][(x+y)]$ જેવી અભિવ્યક્તિ
મંજૂર નથી. તેમ છતાં, $\lambd[y][(x+y)]$ વડે અપાયેલું એકસ્થાની
વિધેય~$f_x(y)$ લઈ શકીએ; તે તેની એકમાત્ર દલીલ~$y$માં~$x$ ઉમેરે
છે. ખરેખર, આ એક વિધેય નથી; દરેક સંખ્યા~$x$ માટે ``$x$ ઉમેરો''
એવાં જુદાં વિધેયોનો પરિવાર છે. હવે~$\lambd[x][f_x]$ તરીકે બીજું
એકસ્થાની વિધેય~$g$ વ્યાખ્યાયિત કરીએ. $x$ પર લાગુ કરતાં~$g(x)$
વિધેય~$f_x$ આપે છે; એટલે એનાં મૂલ્યો પોતે વિધેયો છે. પહેલાં~$g$
ને~$x$ પર અને પછી મળેલા પરિણામને~$y$ પર લાગુ કરીએ તો
$(g(x))y = f_x(y) = x+y$ મળે છે. આ રીતે એકસ્થાની વિધેય~$g$
બે-દલીલવાળા સરવાળાના વિધેય જેવું જ કામ કરે છે. બે-સ્થાની વિધેયને
એકસ્થાની વિધેયમાં (જેના મૂલ્યો એકસ્થાની વિધેયો છે) ફેરવવાની
આ યુક્તિને જ ``કરિંગ'' કહે છે.

હવે~$\lambd$-કલનની વાક્યરચનામાં યોગ્ય એવો દાખલો જોઈએ.
વિધેય $f(x,y) = x$ કેવી રીતે દર્શાવવું? બે દલીલો સ્વીકારીને
પહેલી પરત કરતું વિધેય જોઈએ તો $\lambd[x][\lambd[y][x]]$ લખી
શકીએ. તે ખરેખર દલીલ~$x$ સ્વીકારીને વિધેય~$\lambd[y][x]$
આપતું વિધેય છે. $\lambd[y][x]$ બીજી દલીલ~$y$ સ્વીકારે છે,
પરંતુ તેને અવગણીને હંમેશાં~$x$ જ પરત કરે છે. હવે
$\lambd[x][\lambd[y][x]]$ને બે દલીલો પર લાગુ કરીએ:
\begin{align*}
  (\lambd[x][\lambd[y][x]])MN
  \bredone & (\lambd[y][M])N \\
  \bredone & M
\end{align*}

સામાન્ય રીતે, અમુક પદ~$N$ વડે વ્યાખ્યાયિત $x_1$, \dots,~$x_n$
પ્રાચલોવાળું વિધેય લખવા માટે
$\lambd[x_1][\lambd[x_2][\ldots\lambd[x_n][N]]]$ લખી શકીએ.
તેને~$n$ દલીલો પર લાગુ કરતાં મળે છે:
\begin{multline*}
  (\lambd[x_1][\lambd[x_2][\ldots\lambd[x_n][N]]]) M_1 \dots M_n \bredone\\
  \begin{aligned}
  \bredone {} & (\Subst{(\lambd[x_2][\ldots\lambd[x_n][N]])}{M_1}{x_1}) M_2
  \dots M_n\\
   \eqs {} & (\lambd[x_2][\ldots\lambd[x_n][\Subst{N}{M_1}{x_1}]]) M_2
            \dots M_n \\
  \vdots & \\
  \bredone {} &\Subst{\Subst{N}{M_1}{x_1}\ldots}{M_n}{x_n}
  \end{aligned}
\end{multline*}
છેલ્લી પંક્તિનો સીધો અર્થ છે કે વિધેયની વ્યાખ્યાના શરીરમાં
દરેક~$x_i$ની જગ્યાએ~$M_i$ મૂકવું. અનેક દલીલો વિધેયને આપીએ ત્યારે
આપણને બરાબર એ જ જોઈએ છે.
\sourcecorrection{OLLAM-002}{સ્થિર અંગ્રેજી મૂળની છેલ્લી પંક્તિમાં
\texttt{\textbackslash Subst}ની અંદર~$P$ આવે છે, પરંતુ આ સામાન્ય
વિધેયનું શરીર~$N$ છે અને~$P$ અહીં વ્યાખ્યાયિત જ નથી. તેથી
અહીં~$N$ મૂક્યું છે.}
\end{document}
